% Opens the afternoon deck: the morning in one slide, ending on the advantage,
% which the policy-gradient part builds on.
\section{This Morning}
\begin{frame}{This Morning in One Slide}
\begin{itemize}\itemsep0.7em
    \item \textbf{The problem.} An MDP $(\mathcal{S}, \mathcal{A}, R, P, \gamma)$; a policy
          $\pi$ maps states to actions; the goal is the largest expected discounted return.
    \item \textbf{Values.} $V^{\pi}(s)$: how good a state is. $Q^{\pi}(s,a)$: how good an
          action is in that state. The \textbf{advantage}
          $A^{\pi}(s,a) = Q^{\pi}(s,a) - V^{\pi}(s)$: how much better that action is than
          the policy's average.
    \item \textbf{Q-learning.} Nudge $Q(s,a)$ toward the Bellman target
          $r + \gamma \max_{a'} Q(s',a')$, and act $\varepsilon$-greedily.
\end{itemize}
\vfill
\begin{center}
\itshape This afternoon: skip the values and improve the policy itself.\\
The advantage is how we will decide which actions to make more likely.
\end{center}
\end{frame}
